% Copyright (c) 2026 Eugene Joseph M. Ragasa.
% SPDX-License-Identifier: Apache-2.0
% Author: Eugene Joseph M. Ragasa
% Formal proof completed: 2026-10-19
\documentclass[11pt]{article}
\usepackage[T1]{fontenc}\usepackage{amsmath,amssymb}\usepackage[margin=1in]{geometry}\usepackage[hidelinks]{hyperref}
\title{Proof Dossier: Complex-Argument Harmonic Attainment}\author{Eugene Joseph M. Ragasa}\date{Formal proof completed: 2026-10-19}
\begin{document}\maketitle
\section{Identity and status}
ID \texttt{BG-FOUNDATION-THM-002}; status \texttt{proved}. Lean declaration:
\par\noindent {\scriptsize\path{ProjectKoios.BandGauge.Foundations.secondHarmonicValue_maximizingAngle}}.
\section{Statement and proof}
The angle $\phi=\tfrac12\arg(c+id)$ attains
$b_0+\sqrt{c^2+d^2}$. The polar identities
$\lVert c+id\rVert\cos\arg(c+id)=c$ and
$\lVert c+id\rVert\sin\arg(c+id)=d$ give the result, including $c=d=0$.
\section{Boundary}
No uniqueness is asserted. This is standard supporting analysis. Kernel-reported axioms:
\texttt{propext}, \texttt{Classical.choice}, \texttt{Quot.sound}. No admissions or
scientific axioms. AI-assisted and kernel-checked; no independent expert review recorded.
\end{document}
